\section{The Eastern Common Law, on Its Own Terms}

The 1983 Code binds no Eastern Catholic. Its own first canon says that its canons regard only the Latin Church, and the Eastern Code's first canon says the converse: \latin{Canones huius Codicis omnes et solas Ecclesias orientales catholicas respiciunt}---the canons of this Code regard all and only the Eastern Catholic Churches, unless something else is expressly established as regards relations with the Latin Church. Everything in sections~3 through~6 above therefore stops at that line, and the discipline of twenty-three Churches has to be read from its own book.

\subsection{The witness used here}

There is no Holy See web delivery of the \emph{Codex Canonum Ecclesiarum Orientalium} corresponding to the Latin Code's, and there is no official English version of it at all. This article therefore quotes the Eastern canons from the promulgation text itself: \emph{Acta Apostolicae Sedis} 82 (1990), where the apostolic constitution \emph{Sacri canones} occupies pages 1033--1044 and the Code follows. Each canon quoted below was read at the page image of its \emph{Acta} page in the Holy See's archival PDF of that volume, fetched and hashed on 25 July 2026; the volume's optical-character layer served only to find the pages.\endnote{\emph{Acta Apostolicae Sedis} 82 (1990), archival PDF at the Holy See's \emph{Acta} archive, fetched and hashed 2026-07-25 (SHA-256 recorded in this leaf's research records and in the source-library artifact record). Page images of pp.~1043, 1044, 1061, 1096, 1121, 1142, 1144, 1145, 1213, 1214, 1221, 1222, and 1336 were rendered from that exact PDF and read visually; the OCR layer produced by the volume's scanning software agreed with the visual reading at every phrase quoted here, and where it did not---it drops the entire clause fixing the Code's entry into force---the image governs. One page carrying a canon cited but not quoted in Latin (p.~1220, the opening list of c.~795 §1) was read at the text layer only, and this article marks that lower ceiling here rather than concealing it. The volume PDF is a restricted artifact: the Holy See's portal terms do not grant repository redistribution, so it is identified by hash and retrieval route and its bytes are not reproduced here. English renderings of CCEO canons in this article are the article's own working glosses; the Latin governs. The Canon Law Society of America's Latin--English edition of the Eastern Code exists and was deliberately not used, so that no reading here depends on an unofficial translation.}

The promulgation facts, verified at those pages: \emph{Sacri canones} was given at Rome on 18 October 1990 in the thirteenth year of John Paul II's pontificate, promulgating the Code \latin{pro omnibus Ecclesiis orientalibus catholicis}; and the Code's canons \latin{vim obligandi habere incipiant a die prima mensis Octobris anni MCMXCI, festo Patrocinii Beatae Virginis Mariae in plerisque Orientis Ecclesiis}---began to have binding force on 1 October 1991, the feast of the Protection of the Blessed Virgin Mary in most of the Churches of the East. A code whose \latin{vacatio} is measured to an Eastern feast is announcing something about itself.

\subsection{The interpretive key: canon 2}

That announcement is made explicit in the second preliminary canon, which is the single most important sentence for reading anything that follows: \latin{Canones Codicis, in quibus plerumque ius antiquum Ecclesiarum orientalium recipitur vel accommodatur, praecipue ex illo iure aestimandi sunt}---the canons of the Code, in which the ancient law of the Eastern Churches is for the most part received or adapted, are principally to be assessed from that law.

This is a rule of interpretation with teeth. It instructs the reader not to construe an Eastern canon against the background of the Latin Code but against the background of the Eastern legal tradition the canon is codifying. Applied to the material of this article, it forbids the reading in which the Eastern discipline is a permitted departure from a Latin norm. There is no Latin norm in view. The Second Vatican Council had said the same thing in the language of right rather than of interpretation: \latin{Ecclesias Orientis sicut et Occidentis iure pollere et officio teneri se secundum proprias disciplinas peculiares regendi}---the Churches of the East, as much as those of the West, possess the right and are held by the duty of governing themselves according to their own particular disciplines.\endnote{\emph{Orientalium Ecclesiarum} 5, Latin and English quoted from the Holy See's web deliveries of the Decree, fetched, hashed, and read 2026-07-25. Article~6 of the same Decree adds the corresponding instruction to Easterners themselves---they ``can and should always preserve their legitimate liturgical rite and their established way of life,'' \latin{suos legitimos ritus liturgicos suamque disciplinam semper servare posse et debere}---and to Latins who deal with them. The Decree's protection of Eastern rite and way of life is a distinct claim from \emph{Presbyterorum ordinis} 16's explicit honour for married presbyters; this article keeps them apart.}

\subsection{Canons 373--375: two states, both honoured}

The Eastern Code's treatment of the whole question occupies three canons in the chapter on the rights and obligations of clerics, and the first of them holds both traditions in one sentence:

\begin{studyframe}
\latin{Can.~373 --- Caelibatus clericorum propter regnum coelorum delectus et sacerdotio tam congruus ubique permagni faciendus est, prout fert universae Ecclesiae traditio; item status clericorum matrimonio iunctorum praxi Ecclesiae primaevae et Ecclesiarum orientalium per saecula sancitus in honore habendus est.}

\emph{Working gloss:} The celibacy of clerics, chosen for the sake of the kingdom of heaven and so fitting to the priesthood, is everywhere to be greatly esteemed, as the tradition of the whole Church bears; likewise the state of clerics joined in marriage, sanctioned by the practice of the primitive Church and of the Eastern Churches through the centuries, is to be held in honour.
\end{studyframe}

The canon's architecture is the argument. Two clauses joined by \latin{item}---likewise---each with its own warrant. Celibacy is warranted by \latin{universae Ecclesiae traditio}, the tradition of the whole Church, which is a claim about the East as well as the West; and it is called \latin{sacerdotio tam congruus}, so fitting to the priesthood, in the same comparative register as canon 277 §1's \emph{facilius} and \emph{liberius}. The married state of clerics is warranted by \latin{praxis Ecclesiae primaevae}, the practice of the primitive Church, and by the Eastern Churches' own practice \latin{per saecula}. Neither clause is a concession to the other; neither is described as an exception, a permission, an indult, or a tolerance.

Canon 374 then does what section~2 of this article said the word chastity is for: \latin{Clerici caelibes et coniugati castitatis decore elucere debent; iuris particularis est statuere opportuna media ad hunc finem assequendum adhibenda}---celibate and married clerics alike must shine with the beauty of chastity; it belongs to particular law to establish the suitable means for attaining this end. One obligation for both states, and the concrete means left to the particular law of each Church \latin{sui iuris}. Canon 375 completes the trio by turning to the married cleric's household: in leading family life and in educating their children, married clerics are to give an outstanding example to the rest of the Christian faithful.

What is absent from all three is as instructive as what is present. There is no obligation of perfect and perpetual continence imposed on Eastern clerics by the common law. There is no \latin{ideoque} clause. The Eastern Code does not contain a canon corresponding to CIC canon 277 §1, and the reader who goes looking for one is importing the very assumption canon 2 forbids.

\subsection{Admission of married men: canon 758 §3}

The Eastern Code does not itself decide who may be ordained from among married men. It refers the question outward: \latin{Circa coniugatos ad ordines sacros admittendos servetur ius particulare propriae Ecclesiae sui iuris vel normae speciales a Sede Apostolica statutae}---as regards married men to be admitted to sacred orders, the particular law of one's own Church \latin{sui iuris} is to be observed, or special norms established by the Apostolic See.

Three things follow, and the third is where most confusion in this area is generated. First, the common law of the Eastern Churches contains no impediment corresponding to CIC canon 1042 1°. The list at CCEO canon 762 §1 runs to eight numbers---psychic illness, the delict of apostasy, heresy or schism, attempted marriage by one impeded (3°), voluntary homicide or procured abortion, grave self-mutilation or attempted suicide, the reserved act of orders, a forbidden office or administration, and the neophyte---and the married man appears nowhere in it.\endnote{CCEO c.~762 §1, read at the page image of \emph{Acta Apostolicae Sedis} 82 (1990) 1214. The same page carries two further comparative details worth noting. Canon 760 §2 provides the Eastern law's route to a stable diaconate---a candidate not destined for the priesthood may be ordained deacon after the third year of the studies of c.~354---and canon 761's required declaration is verbally fuller than CIC c.~1036: the Eastern candidate testifies that he will receive the sacred order \emph{and the obligations annexed to that order}, \latin{obligationes eidem ordini adnexas}, which in the East is the natural formula precisely because those obligations differ with the candidate's state.} Second, the competent legislator is ordinarily the Church \latin{sui iuris} itself, through its own particular law---which means that the practice of one Eastern Church establishes nothing whatever about another. Third, the alternative in the canon, \latin{normae speciales a Sede Apostolica statutae}, is a standing opening for the Apostolic See to legislate specially in this matter, and it has been used, most consequentially in the twentieth century for the Eastern diasporas.

That last opening has been used, and heavily. The restrictions imposed on married Eastern clergy in the Western diasporas across the twentieth century, and the \emph{Pontificia Praecepta de clero uxorato orientali} published in \emph{Acta Apostolicae Sedis} 106 (2014) 496--499, are acts of exactly the kind canon 758 §3 contemplates. Section~8 works through them at the \emph{Acta} text. The point to carry here is structural: because canon 758 §3 names the Apostolic See's special norms as a source alongside each Church's own particular law, a restriction imposed by such norms is not a departure from Eastern law but an operation of it---and so is its relaxation.

\subsection{Canon 804: the point of complete agreement}

\latin{Invalide matrimonium attentat, qui in ordine sacro est constitutus}---he who is constituted in a sacred order invalidly attempts marriage. Word for word the substance of CIC canon 1087, and the two codes are here identical: in neither Church may a man marry after ordination. What the East permits is the ordination of a man already married. What neither permits is marriage afterwards.

The corollaries follow the Latin pattern exactly. The impediment is of ecclesiastical law; its dispensation is reserved to the Apostolic See (c.~795 §1 1° with §2); and in danger of death the local Hierarch may dispense from every ecclesiastical impediment \latin{excepto impedimento ordinis sacri sacerdotii}---except the impediment of the sacred order of priesthood (c.~796 §1), the same narrow reservation the Latin Code makes at canon 1079 §1. Two codes, drafted separately, reserving the same thing in the same words at the same point.

\subsection{The episcopate, exactly as the canon states it}

Eastern bishops are, in fact, unmarried men, and the reason is often given as a rule of celibacy. The canon says something more precise. Among the requirements for a candidate to be held suitable for the episcopate, canon 180 3° lists: \latin{vinculo matrimonii non ligatus}---not bound by the bond of marriage.

That is a requirement about a bond, not about a history. A widower is not bound by the bond of marriage, and the canon does not exclude him; nor does it use the word \latin{caelebs}, which the Code has and uses elsewhere (cc.~373, 374, 376). Paul~VI's observation that ``in the East only celibate priests are ordained bishops'' describes the practice accurately, and the practice---bishops drawn from the monastic clergy---is older and stronger than the canon. But the enacted requirement is the narrower one, and an article about what the law says should say what the law says.\endnote{CCEO c.~180, read at the page image of \emph{Acta Apostolicae Sedis} 82 (1990) 1096. \emph{Sacerdotalis caelibatus} n.~40, quoted from the Holy See's English web delivery of the encyclical, fetched, hashed, and read 2026-07-25. The practice by which Eastern bishops are taken from among monastic or celibate clergy is not established here from particular law; the claim made is only that the common-law requirement is stated as absence of a marriage bond.}

\subsection{A code written for a married clergy}

The clearest evidence that the Eastern Code treats married clerics as ordinary rather than exceptional is not in the canons about marriage at all. It is in the ones about money and health.

Canon 390 §1 gives clerics the right to suitable support and just remuneration, and adds: \latin{si agitur de clericis coniugatis, consulere debet etiam eorum familiae sustentandae}---if it concerns married clerics, it must also provide for the support of their family. Section 2 extends the right to social security and health assistance ``for themselves and their family, if they are married.'' Canon 285 §2 provides that if a presbyter to be named as pastor is joined in marriage, good morals are required in his wife and in the children living with him. Canon 397 assumes, in setting out the Patriarch's competence to grant loss of the clerical state, that some clerics simply are not bound by any obligation of celibacy: he may grant it, with the consent of the Synod of Bishops of the patriarchal Church, to clerics domiciled within its territory \latin{qui obligatione caelibatus non tenentur aut, si tenentur, dispensationem ab hac obligatione non petunt}---who are not bound by the obligation of celibacy, or who, if they are bound, do not seek dispensation from it; in other cases the matter goes to the Apostolic See.\endnote{CCEO cc.~285 §2, 390 §§1--2, 394--398, read at the page images of \emph{Acta Apostolicae Sedis} 82 (1990) 1121, 1144, and 1145. Canon 394 is the Eastern parallel to CIC c.~290 and repeats its opening premise verbatim in substance---\latin{Sacra ordinatio semel valide suscepta numquam irrita fit}---while naming the second mode of loss as the penalty of \emph{deposition} rather than dismissal. The division of labour in c.~397 is the point of interest here: the Patriarch's competence is drawn precisely along the line of the celibacy obligation, because in the East that line runs through the clergy rather than around it.}

A body of law that provides pensions for clergy wives and screens the morals of a pastor's children is not making an allowance. It is legislating for its ordinary case.

\subsection{The two codes side by side}

\begin{threestable}{Question}{Latin Church (CIC 1983)}{Eastern Catholic Churches (CCEO 1990)}
Is there a general obligation of continence on clerics? & Yes: c.~277 §1 obliges perfect and perpetual continence, and therefore celibacy. & No such common-law obligation. Cc.~373--375 esteem celibacy, honour the married state, and oblige chastity in both (c.~374).\\
May a married man be ordained? & Deacon: yes if legitimately destined to the permanent diaconate (c.~1042 1°), at 35 with his wife's consent (c.~1031 §2). Presbyter: only by dispensation from c.~1042 1° reserved to the Apostolic See (c.~1047 §2 3°). & Yes in principle; the question is referred to the particular law of each Church \latin{sui iuris} or to special norms of the Apostolic See (c.~758 §3). No impediment attaches to the married state as such (c.~762 §1).\\
May a cleric marry after ordination? & No: those in sacred orders invalidly attempt marriage (c.~1087). & No: identical rule at c.~804.\\
Who dispenses that impediment? & The Apostolic See (c.~1078 §2 1°); in danger of death, the local ordinary except for the presbyterate (c.~1079 §1). & The Apostolic See (c.~795 §1 1° and §2); in danger of death, the local Hierarch except for the priesthood (c.~796 §1).\\
Must a bishop be unmarried? & The question does not arise separately: all Latin candidates are already bound by c.~277 §1. & Yes, stated as a requirement of suitability: not bound by the bond of marriage (c.~180 3°).\\
How is the obligation of celibacy undertaken? & Publicly, in the prescribed rite, before the diaconate (c.~1037). & Not by a common-law act; where the obligation exists it arises from the particular law or the candidate's state.\\
Does loss of the clerical state release from celibacy? & No; only the Roman Pontiff dispenses (c.~291). & No; only the Roman Pontiff dispenses (c.~396). The Patriarch may grant loss of the clerical state itself in defined cases (c.~397).\\
Penalty for attempted marriage? & \latin{Latae sententiae} suspension, escalating to dismissal after warning (c.~1394 §1). & Deposition, and not \latin{latae sententiae}: \latin{Clericus, qui prohibitum matrimonium attentavit, deponatur} (c.~1453 §2).\\
Penalty for scandalous unchastity? & Suspension, escalating after warning (c.~1395 §1), for a cleric persisting in an external sin \latin{contra sextum Decalogi praeceptum}. & Suspension, escalating to deposition (c.~1453 §1), for a cleric persisting in an external sin \latin{contra castitatem}.\\
\end{threestable}

Read down the third column and the shape of the Eastern law appears: it agrees with the Latin Code precisely where the ancient discipline of the whole Church agreed---no marriage after orders, an unmarried episcopate, dispensation reserved---and it diverges precisely where the Latin Church added a determination of her own. The divergence is not the East's departure. It is the West's development, and the East's continuity.
